\clearpage
\section*{Checking what is read from eLMS}
\textit{September 30, 2026. Agent-drafted work on branch \texttt{explore-alderman-measures}; Jacob asked that the
extraction be checked carefully before more measures are built on it.}

Three readings of the eLMS records were audited: the application forms, the sites, and the districts before and after.
The form parser's plausibility flags fell from 195 records to 174 after nine rule fixes, most of the rest large
planned developments whose values are right, and both hand-read holdouts improved without new errors
(\path{tasks/audits/zoning_record_validation}, sections 3 and 8). Applications without a ward fell from 157 to 26 by
reading the title address again onto the 2013 street centerlines where the geocoder failed; where both place a site,
the reread lies a median of 52 feet from the geocoder's point. The geocoder's responses are now a preserved snapshot
(\path{tasks/assign_zoning_amendment_wards}), so that re-reading the ordinances does not re-query a live service.

The districts were the weakest reading: the kind of change was unknown for 19 percent of eLMS applications. The
ordinance sentence was missed when it was worded or misread in ways the reader did not allow (``changing the'' without
``all'', ``District, as shown on Map'' without ``symbols'', ``10 those of'', ``vvhich'' for the word ending the new
district), when OCR misread a code (``83-2'', ``RMS'' for RM-5), when a substitute's file name began ``S0'', and when
a planned development's statements were headed ``Plan of Development Statements''. With these read
(\path{tasks/clean_zoning_map_amendments}), the kind is unknown for 4 percent of applications and 16 percent of
aldermen's amendments; the largest group of the rest changes a planned development, which has no floor-area ratio, to
a district. Two independent checks agree. The City's zoning map names the ordinance record that set each district
since 2011; for 98 percent of the 4,105 passed amendments on it, the map has the district read, newly read districts
as often as the others. Against the districts the Journals print for 2010--2011 passages, agreement rose from 94 to 97
percent while the number compared rose from 215 to 271. Where an introduced ordinance and its substitute name
different districts (99 passed amendments), the map has the substitute's for 85, which supports reading the
substitute.

Stalls still carry no alderman signal, but several results of the September 29 entry change. Reilly's passages took 5
percent longer than those of the same kind before and 23 percent shorter now: 97 of his 116 applications are planned
developments, 51 of which had been of unknown kind, and among planned developments his are fast, as the matched
comparison of large projects also finds (0.69 log points below their matches). Conway's large downtown projects take
longer than their matches by 0.80 log points (standard error 0.31), and by 0.39 without 30 N LaSalle and 37 S
Sangamon. In the production measures (\path{tasks/estimate_alderman_zoning_measures}) days to passage are now shrunk
to zero for every alderman, because the kind of change explains more of the differences and the true variance is
estimated from all aldermen, those with few passages included. And aldermen's own downzonings in eLMS had been
undercounted: batches such as Maldonado's 19 of December 2018 (16 of them RS-3 to RS-1) and Scott's 16 of April 2022
(to B2-1) were of unknown kind. The term of 2019--2023 had 120, not 46, and the fall from 540 in 2003--2007 is less
steep. The turnover tests of new construction are unchanged (\path{tasks/audits/alderman_turnover_volumes}).

\textit{Open.} \path{tasks/explore_alderman_measures} reads the production measures and was not rebuilt. An ordinance
changing several separate areas is read as its first area's district before and its last area's after.
